\documentclass[leqno]{article}
% Shared mark scheme preamble for resources.danieljsmith.org: the pack
% preamble, unchanged, then the scheme's own layer. A scheme's .tex is
% rendered from its YAML (scripts/README.md, Mark schemes), never edited.
\input{\CurrentFilePath/pack_preamble.tex}

% The page: equal side margins, so the table runs the full width of the
% header rule. The footers' right edge follows the text's.
\geometry{left=0.8in, right=0.8in}
\makeatletter
\renewcommand{\djs@footR}[1]{%
  \put(\dimexpr1in+\hoffset+\oddsidemargin+\textwidth\relax,
    -\dimexpr\paperheight-0.32in\relax){\makebox[0pt][r]{#1}}}
\makeatother

\usepackage[export]{adjustbox}
\usepackage{tabularray}
\definecolor{msGrid}{RGB}{176, 172, 164}
\definecolor{msTint}{RGB}{245, 243, 239}

% The header on every page, the front page included, and the booklet label
% under the contents-page title. \djsQuestionPageStyle stays empty.
\newcommand{\djsSchemeHeader}[1]{%
  \djsPackHeader[markscheme]{\textit{Mark scheme} - #1}%
  \renewcommand{\djsPackTitle}{#1}%
  \renewcommand{\djsBookletLabel}{Mark scheme}%
}

% The front page: the contents heading, spaced as tocloft spaces the packs'
% "Questions", then the contents table. A question's name and page number
% both link to its first page.
\newcommand{\msContentsHeading}{%
  \vskip\cftbeforetoctitleskip
  {\noindent\cfttoctitlefont\contentsname\par}%
  \vskip\cftaftertoctitleskip}
\NewTblrEnviron{mscontents}
\SetTblrInner[mscontents]{
  colspec = {Q[wd=26mm, l] X[l] Q[wd=16mm, c] Q[wd=14mm, r]},
  width = \linewidth,
  rowsep = 4pt, colsep = 6pt,
  hlines = {0.4pt, msGrid}, vlines = {0.4pt, msGrid},
  hline{2} = {2.2pt, scarlet},
  row{1} = {font=\small\bfseries},
  row{Z} = {font=\bfseries},
  column{2} = {font=\small},
}
\newcommand{\msContentsQuestion}[1]{\hyperref[q:#1]{Question #1}}
\newcommand{\msContentsPage}[1]{\hyperref[q:#1]{\pageref*{q:#1}}}
\newcommand{\msSep}{\hspace{0.9em}}

% The conventions pages' layout: headings in the documents' own bold, tables
% ruled and padded as the scheme's, roomy lists.
\newcommand{\msHead}[1]{%
  \par\vspace{18pt}%
  {\noindent\bfseries #1\par}%
  \vspace{6pt}}
\NewTblrEnviron{mstblr}
\SetTblrInner[mstblr]{rowsep = 6pt, colsep = 6pt,
  hlines = {0.4pt, msGrid}, vlines = {0.4pt, msGrid}}
\newenvironment{msList}
  {\begin{itemize}[leftmargin=1.2em, itemsep=5pt, topsep=3pt, parsep=0pt]}
  {\end{itemize}}

% A question's first page: the link target and the PDF bookmark, as
% \questionitem gives the packs.
\newcommand{\msQnum}{}
\newcommand{\msQuestion}[1]{%
  \clearpage
  \renewcommand{\msQnum}{#1}%
  \phantomsection\label{q:#1}%
  \addcontentsline{toc}{subsection}{\protect\djsTocEntry{#1}}}

% The scheme's table: Question 16mm, Scheme the rest, Marks 16mm, Notes
% 44mm; every cell ruled in msGrid; the scarlet rule under the head row. A
% long table has no caption or "continued" foot; a continuation page opens
% with a muted line above the repeated head row. The rules across a part's
% rows, the alternative's tint and the total row are the rendered table's.
\NewTblrTheme{ms}{%
  \DefTblrTemplate{firsthead}{default}{}%
  \DefTblrTemplate{middlehead,lasthead}{default}{%
    {\small\itshape\color{djsMuted}Question \msQnum\ continued\par}%
    \vspace{6pt}}%
  \DefTblrTemplate{firstfoot,middlefoot,lastfoot}{default}{}%
}
\NewTblrEnviron{msscheme}
\SetTblrOuter[msscheme]{long, theme=ms, entry=none, label=none}
\SetTblrInner[msscheme]{
  colspec = {Q[wd=16mm, l] X[l] Q[wd=16mm, c] Q[wd=44mm, l]},
  width = \linewidth,
  rowsep = 6pt, colsep = 6pt,
  rowhead = 1,
  vlines = {0.4pt, msGrid},
  hline{1} = {0.4pt, msGrid}, hline{2} = {2.2pt, scarlet},
  hline{Z} = {0.4pt, msGrid},
  column{1} = {font=\bfseries}, column{3} = {font=\bfseries},
  column{4} = {font=\small},
  row{1} = {rowsep=4pt, font=\small\bfseries},
}

% Two marks on a row: each note starts with its code.
\newcommand{\msNote}[2]{\textbf{#1}\enspace #2}
% A full-row figure: the worked solution's own picture, scaled uniformly and
% only down, to the table less the Question column. The picture is drawn
% with the line width it has in the solutions document, at its list depth
% there (the optional argument, 1 by default: a part's working), so one
% sized by \linewidth keeps its size; tabularray sets the cell's \linewidth
% to the cell's measured width, which would shrink it at every pass. The
% widths are measured once, in a solution box at the packs' text width
% (pack_preamble.tex's margins, 0.8in and 1in).
\newlength{\msSpanW}
\makeatletter
\newcommand{\ms@keepLW}[1]{%
  \expandafter\xdef\csname ms@solnLW@#1\endcsname{\the\linewidth}}
\AtBeginDocument{%
  \setlength{\msSpanW}{\dimexpr\textwidth-16mm-28pt\relax}%
  \setbox\z@\vbox{%
    \begin{minipage}{\dimexpr\paperwidth-1.8in\relax}
      \begin{djsSolution}
        \ms@keepLW{0}%
        \begin{enumerate}\item \ms@keepLW{1}%
        \begin{enumerate}\item \ms@keepLW{2}%
        \begin{enumerate}\item \ms@keepLW{3}%
        \end{enumerate}\end{enumerate}\end{enumerate}
      \end{djsSolution}
    \end{minipage}}}
\newcommand{\msFigure}[2][1]{%
  \adjustbox{max width=\msSpanW, max totalheight=0.6\textheight, valign=t}{%
    \hsize=\csname ms@solnLW@#1\endcsname
    \linewidth=\hsize
    #2}}
\makeatother
% A figure mark's feature, in the Scheme column.
\newcommand{\msFeature}[1]{\textit{#1}}
% An alternative's title row.
\newcommand{\msAlternative}[1]{\textit{Alternative:} #1}
% Drafts only: a queried mark's code, the query's why under its note, and a
% flag in its part's place.
\newcommand{\msQueried}[2]{#1~\(\to\) #2?}
\newcommand{\msQuery}[1]{{\itshape\color{djsMuted}Query: #1}}
\newcommand{\msFlag}[2]{\textbf{#1:}\enspace #2}


\djsSchemeHeader{Matrix Determinants \& Inverses}
\djsSitePath{Core Pure}{Vectors, Matrices \& Linear Transformations}

\djsLastUpdated{26 October 2026}
\begin{document}
\djsSetQuestionCount{10}
\djsAddMarks{1}{5}
\djsAddMarks{2}{6}
\djsAddMarks{3}{10}
\djsAddMarks{4}{7}
\djsAddMarks{5}{5}
\djsAddMarks{6}{8}
\djsAddMarks{7}{6}
\djsAddMarks{8}{8}
\djsAddMarks{9}{14}
\djsAddMarks{10}{9}
\vspace*{6pt}
\djsContentsTitle
\vspace{30pt}
\msContentsHeading
\begin{mscontents}{}
Question & Marks by part & Marks & Page \\
\msContentsQuestion{1} & (a)\,2\msSep (b)\,3 & 5 & \msContentsPage{1} \\
\msContentsQuestion{2} & (a)\,2\msSep (b)\,4 & 6 & \msContentsPage{2} \\
\msContentsQuestion{3} & (a)(i)\,2\msSep (a)(ii)\,4\msSep (b)\,4 & 10 & \msContentsPage{3} \\
\msContentsQuestion{4} & & 7 & \msContentsPage{4} \\
\msContentsQuestion{5} & (a)\,3\msSep (b)\,2 & 5 & \msContentsPage{5} \\
\msContentsQuestion{6} & & 8 & \msContentsPage{6} \\
\msContentsQuestion{7} & (a)\,2\msSep (b)\,4 & 6 & \msContentsPage{7} \\
\msContentsQuestion{8} & (a)\,2\msSep (b)\,2\msSep (c)\,1\msSep (d)\,3 & 8 & \msContentsPage{8} \\
\msContentsQuestion{9} & (a)\,2\msSep (b)(i)\,6\msSep (b)(ii)\,3\msSep (c)(i)\,1\msSep (c)(ii)\,2 & 14 & \msContentsPage{9} \\
\msContentsQuestion{10} & (a)\,4\msSep (b)\,5 & 9 & \msContentsPage{10} \\
\SetCell[c=2]{r} Total & & 78 & \\
\end{mscontents}
\newpage
% The conventions pages, shared by every pack's mark scheme: one source of
% truth, \input after the front page, with no strand or pack logic. The keys
% on the first side, the rules on the second. The scheme layer defines
% \msHead, msList and mstblr.
\vspace*{6pt}%
{\noindent\cfttoctitlefont Reading the mark scheme\par}%
\vspace{10pt}%
These rules apply throughout, unless a note says otherwise.

\msHead{Marks}
Every mark is worth one, so \textbf{M1 A1} on a row is two marks. Only marks
coded ft follow through.\par
\vspace{10pt}%
\begin{mstblr}{
  colspec = {Q[wd=15mm, l] X[l]},
  width = \linewidth,
  hline{2} = {2.2pt, scarlet},
  column{1} = {font=\bfseries},
  row{1} = {font=\small\bfseries, rowsep=4pt},
}
Code & Meaning \\
M1 & Method: a correct method, applied to this question. A slip doesn't cost it if the method is clear. A formula must be used, not just quoted. \\
dM1 & Dependent method: as M1, but only if the method mark before it was earned, unless its note names another. \\
A1 & Accuracy: a correct answer or step. It needs the method mark it follows. \\
A1ft & As A1, following through from an earlier answer that the note names. \\
A1~AG & The last mark of a part whose answer is printed in the question: complete working, with no errors. \\
B1 & A correct value, fact or statement. It needs no method mark. \\
B1ft & As B1, following through from an earlier answer that the note names. \\
E1 & A written explanation, interpretation or conclusion. Concluding a proof by induction needs every earlier mark. \\
\end{mstblr}

\msHead{Abbreviations}
ft and AG qualify a code in the Marks column. The others appear in the
notes.\par
\vspace{10pt}%
\begin{mstblr}{
  colspec = {Q[wd=17mm, l] X[l] Q[wd=15mm, l] X[l]},
  width = \linewidth,
  column{1} = {font=\bfseries}, column{3} = {font=\bfseries},
}
ft & follow through & AG & answer given in the question \\
cao & correct answer only, with no follow-through & cso & correct solution only: no errors in the part \\
oe & or an equivalent form & awrt & any answer that rounds to this value \\
soi & seen, or implied by later working & isw & ignore working after a correct answer \\
SC & special case: credit outside the main method & their & the student's own earlier value \\
condone & a slip that doesn't cost the mark & [\;\,] & working the mark doesn't need \\
s.f. & significant figures & d.p. & decimal places \\
\end{mstblr}

\newpage
\msHead{Reading the scheme}
\begin{msList}
\item Each row is a line of working, with the marks it earns beside it.
  Where a row carries two marks, each note starts with its code.
\item \textbf{[\textit{n}]} ends each part with its total.
\item A shaded \textit{Alternative} is a second method for the same total.
  Mark an answer by the method it follows.
\item Any other valid method earns the equivalent marks, unless the question
  names a method or says \textit{hence}.
\item A sketch earns one B1 for each feature named below its figure.
\end{msList}

\msHead{Accuracy and answers}
\begin{msList}
\item Unless a question says otherwise, give answers to 3 significant
  figures, and angles in degrees to 1 decimal place. Angles in radians take
  3 significant figures.
\item A final answer with more figures is accepted if it rounds to the
  scheme's value; one with fewer is not. An answer spoiled by rounding too
  early loses its accuracy mark.
\item Where a question asks for an exact answer, a decimal loses the final
  accuracy mark. Otherwise either is accepted.
\item Wrong or missing units cost nothing unless the question asks for them.
\item A correct answer with no working earns full marks, unless the question
  needs the working: show that, prove, hence, a method it names, or an exact
  answer. A wrong answer with no working earns nothing. Calculator answers
  are treated the same way.
\item A correct answer, once reached, isn't spoiled by later working (isw).
\end{msList}

\msHead{Value of $\boldsymbol{g}$}
\begin{msList}
\item Take $g = 9.8\,\mathrm{m\,s^{-2}}$ unless a question states a value,
  which is then the only one accepted. Answers that use $g = 9.8$ may be
  given to 2 or 3 significant figures.
\item Where a question doesn't state $g$, an answer worked with $g = 10$
  throughout earns the same marks, keeping any value printed in the
  question. The worked solutions list the final answers with $g = 10$.
\end{msList}


\msQuestion{1}
\begin{msscheme}{
  hline{3-4,6-8} = {2-Z}{0.4pt, msGrid},
  hline{5,9} = {0.4pt, msGrid},
  row{9} = {font=\bfseries},
}
Question & Scheme & Marks & Notes \\
1(a) & {\(\det\mathbf{B} = 3 - a\)\\[4pt] \(\mathbf{B}^{-1} = \frac{1}{\ldots}\begin{pmatrix}3 & -a\\ -1 & 1\end{pmatrix}\)} & M1 & {Uses either the swap of the leading diagonal with the other entries negated, or division by a determinant of \(\pm(3 - a)\).} \\
& {\(\mathbf{B}^{-1} = \frac{1}{3-a}\begin{pmatrix}3 & -a\\ -1 & 1\end{pmatrix}\)} & A1 & {oe, such as the factor taken into the entries.} \\
\nopagebreak & & [2] & \\
1(b) & {\(2\mathbf{I} - \frac12\mathbf{B} = \begin{pmatrix}\frac32 & -\frac{a}{2}\\ -\frac12 & \frac12\end{pmatrix}\)} & B1 & {oe, such as \(\frac12\begin{pmatrix}3 & -a\\ -1 & 1\end{pmatrix}\).} \\
& {Equates with their \(\mathbf{B}^{-1}\):\\[4pt] \(\frac{1}{3-a} = \frac12 \Rightarrow a = \ldots\)} & M1 & {Equates their \(\mathbf{B}^{-1}\) in terms of \(a\) with their \(2\mathbf{I} - \frac12\mathbf{B}\), and solves an equation in \(a\) from the scalar factors or one entry.} \\
& {\(a = 1\)} & A1 & {\(a = 1\) only, with any other value from comparing entries rejected.} \\
\nopagebreak & & [3] & \\
\nopagebreak \SetCell[c=2]{r} Total for Question 1 & & 5 & \\
\end{msscheme}

\msQuestion{2}
\begin{msscheme}{
  hline{3-4,6-9} = {2-Z}{0.4pt, msGrid},
  hline{5,10} = {0.4pt, msGrid},
  row{10} = {font=\bfseries},
}
Question & Scheme & Marks & Notes \\
2(a) & {\(\mathbf{BA} = \begin{pmatrix}(k+1)+(1-2k)+(2k-1) & \ldots\\ \ldots & \ldots\end{pmatrix}\)} & M1 & {Forms a \(2\times 2\) matrix by multiplying rows of \(\mathbf{B}\) by columns of \(\mathbf{A}\), with at least two entries correct. Need not be simplified.} \\
& {\(\mathbf{BA} = \begin{pmatrix}k+1 & 2\\ 3k & 6-k\end{pmatrix}\)} & A1 & {All four entries correct and simplified.} \\
\nopagebreak & & [2] & \\
2(b) & {\((k+1)(6-k) - 2(3k) = 0\)} & M1 & {Sets \(ad - bc\) of their \(2\times 2\) \(\mathbf{BA}\) equal to zero, in either order of subtraction. Condone sign slips in expanding.} \\
& {\(k^2 + k - 6 = 0\)} & A1 & {oe three-term quadratic.} \\
& {\((k+3)(k-2) = 0 \Rightarrow k = \ldots\)} & dM1 & {Solves their three-term quadratic by factorising, the formula or completing the square, leading to \(k = \ldots\); a calculator is allowed.} \\
& {\(k = -3\), \(k = 2\)} & A1 & {Both values and no others.} \\
\nopagebreak & & [4] & \\
\nopagebreak \SetCell[c=2]{r} Total for Question 2 & & 6 & \\
\end{msscheme}

\msQuestion{3}
\begin{msscheme}{
  hline{3-4,6-9,11-13} = {2-Z}{0.4pt, msGrid},
  hline{5,10,14} = {0.4pt, msGrid},
  row{14} = {font=\bfseries},
}
Question & Scheme & Marks & Notes \\
3(a)(i) & {\(C\), \(E\), \(F\) for Commuter, Electric, Folding} & B1 & \\
& {each the number of that model manufactured in 2023} & B1 & {Needs both the number of bicycles and the year 2023; one statement may cover all three.} \\
\nopagebreak & & [2] & \\
3(a)(ii) & {\(C + E + F = 1000\)} & B1 & \\
& {\(F = C + E - 200\)} & B1 & {oe, such as \(C + E - F = 200\).} \\
& {\(1.10C + 0.95E + 1.05F = 1000 \times 1.038\)} & M1 & {Applies a percentage multiplier to each variable and equates the sum to a 2024 total, or the sum of the changes to the increase. Condone one wrong multiplier.} \\
& {\(1.10C + 0.95E + 1.05F = 1038\)} & A1 & {oe, such as \(0.1C - 0.05E + 0.05F = 38\).} \\
\nopagebreak & & [4] & \\
3(b) & {\(\begin{pmatrix}1&1&1\\-1&-1&1\\1.10&0.95&1.05\end{pmatrix}\begin{pmatrix}C\\E\\F\end{pmatrix} = \begin{pmatrix}1000\\-200\\1038\end{pmatrix}\)} & M1 A1ft & {\msNote{M1}{Writes their equations from (a)(ii) as a \(3 \times 3\) coefficient matrix times \(\begin{pmatrix}C\\E\\F\end{pmatrix}\) equal to a column of constants, or a \(2 \times 2\) one after eliminating a variable. Without a matrix equation, (b) scores 0.}\\ \msNote{A1ft}{Correct matrix equation for their (a)(ii) equations; any consistent order of rows and variables.}} \\
& {\(\begin{pmatrix}C\\E\\F\end{pmatrix} = \their{\begin{pmatrix}1&1&1\\-1&-1&1\\1.10&0.95&1.05\end{pmatrix}}^{-1}\begin{pmatrix}1000\\-200\\1038\end{pmatrix}\)\\[4pt] \(\phantom{\begin{pmatrix}C\\E\\F\end{pmatrix}} = \ldots\)} & dM1 & {Multiplies the constant column on the left by the inverse of their coefficient matrix and obtains values. The inverse may be found by calculator and need not be shown. Not for elimination alone.} \\
& {\(C = 320,\ E = 280,\ F = 400\)} & A1 & {cao. Bare correct values after a correct matrix equation imply the dM1; values found by elimination alone, with no inverse used, score dM0 A0.} \\
\nopagebreak & & [4] & \\
\nopagebreak \SetCell[c=2]{r} Total for Question 3 & & 10 & \\
\end{msscheme}

\msQuestion{4}
\begin{msscheme}{
  hline{3-8} = {2-Z}{0.4pt, msGrid},
  hline{9} = {0.4pt, msGrid},
  row{9} = {font=\bfseries},
}
Question & Scheme & Marks & Notes \\
4 & {\(A + D + P = 580\)} & B1 & \\
& {\(A - D - P = 20\)} & B1 & {oe, such as \(A = D + P + 20\).} \\
& {\(0.95A + 1.10D + 1.15P = \ldots\)} & M1 & {Applies a percentage multiplier to each of \(A\), \(D\) and \(P\) and equates the sum to a new total, or the sum of the changes to the increase. Condone one wrong multiplier.} \\
& {\(0.95A + 1.10D + 1.15P = 600\)} & A1 & {oe, such as \(19A + 22D + 23P = 12\,000\) or \(-0.05A + 0.1D + 0.15P = 20\).} \\
& {\(\begin{pmatrix}1&1&1\\1&-1&-1\\19&22&23\end{pmatrix}\begin{pmatrix}A\\D\\P\end{pmatrix} = \begin{pmatrix}580\\20\\12\,000\end{pmatrix}\)} & M1 & {Writes their three equations as a \(3 \times 3\) coefficient matrix times \(\begin{pmatrix}A\\D\\P\end{pmatrix}\) equal to a constant vector, or a \(2 \times 2\) one after eliminating a variable. Without a matrix equation, only B1 B1 M1 A1.} \\
& {\(\begin{pmatrix}A\\D\\P\end{pmatrix} = \begin{pmatrix}1&1&1\\1&-1&-1\\19&22&23\end{pmatrix}^{-1}\begin{pmatrix}580\\20\\12\,000\end{pmatrix} = \ldots\)} & dM1 & {Multiplies their constant vector on the left by the inverse of their coefficient matrix, by calculator or otherwise, to obtain values. Not for elimination alone.} \\
& {\(A = 300,\ D = 140,\ P = 140\)} & A1 & {cao. Bare correct values after a correct matrix equation imply the dM1; values found by elimination alone, with no inverse used, score dM0 A0.} \\
\nopagebreak \SetCell[c=2]{r} Total for Question 4 & & 7 & \\
\end{msscheme}

\msQuestion{5}
\begin{msscheme}{
  hline{3-5,7-8} = {2-Z}{0.4pt, msGrid},
  hline{6,9} = {0.4pt, msGrid},
  row{9} = {font=\bfseries},
}
Question & Scheme & Marks & Notes \\
5(a) & {\(\det\mathbf{A} = 2(0+2) - k(1-2k) + 1(-1-0)\)} & M1 & {Expands along any row or column, or by the rule of Sarrus, with all three terms. Condone sign and arithmetic slips.} \\
& {\(\det\mathbf{A} = 2k^2 - k + 3\)} & A1 & \\
& {\(b^2 - 4ac = (-1)^2 - 4(2)(3) = -23 < 0\)\\[4pt] \(\Rightarrow \det\mathbf{A} \neq 0\) for all real \(k\)\\[4pt] \(\Rightarrow \mathbf{A}\) is non-singular} & A1~AG & {Correct discriminant, or \(2\left(k-\frac14\right)^2+\frac{23}{8} > 0\), with \(\det\mathbf{A} \neq 0\) or non-singular concluded for all \(k\).} \\
\nopagebreak & & [3] & \\
5(b) & {\(2\times\tfrac12 + k\times\tfrac14 + 1\times\left(-\tfrac14\right) = 1\)\\[4pt] \(\Rightarrow k = \ldots\)} & M1 & {Equates an entry of \(\mathbf{AA}^{-1}\) or \(\mathbf{A}^{-1}\mathbf{A}\) involving \(k\) to that of \(\mathbf{I}\), or uses another valid link, and solves for \(k\). Condone arithmetic slips.} \\
& {\(k = 1\)} & A1 & {\(k = 1\) only; A0 if another value found is not rejected.} \\
\nopagebreak & & [2] & \\
\nopagebreak \SetCell[c=2]{r} Total for Question 5 & & 5 & \\
\end{msscheme}

\msQuestion{6}
\begin{msscheme}{
  hline{3-7} = {2-Z}{0.4pt, msGrid},
  hline{8} = {0.4pt, msGrid},
  row{8} = {font=\bfseries},
}
Question & Scheme & Marks & Notes \\
6 & {\(x + y + z = 900\,000\)} & B1 & \\
& {\(1.25x + 0.8y + 1.5z = 1\,050\,000\)} & M1 A1 & {\msNote{M1}{Applies a percentage multiplier to at least two of the 2023 variables and equates the sum to \(1\,050\,000\), or the sum of the changes to \(150\,000\).}\\ \msNote{A1}{oe, such as \(25x + 16y + 30z = 21\,000\,000\) or \(0.25x - 0.2y + 0.5z = 150\,000\).}} \\
& {\(16(1.25x) + 15(0.8y) + 12(1.5z) = 15\,120\,000\)} & M1 A1 & {\msNote{M1}{Multiplies each profit per pass by their 2024 sales of that type and equates the sum to the total profit. Condone multiplier slips.}\\ \msNote{A1}{oe, such as \(20x + 12y + 18z = 15\,120\,000\), or consistent working in millions.}} \\
& {\(\begin{pmatrix}1&1&1\\25&16&30\\20&12&18\end{pmatrix}\begin{pmatrix}x\\y\\z\end{pmatrix} = \begin{pmatrix}900\,000\\21\,000\,000\\15\,120\,000\end{pmatrix}\)} & M1 & {Writes their three equations as a \(3 \times 3\) coefficient matrix times \((x, y, z)\) equal to their constants, or a \(2 \times 2\) one after eliminating a variable. Without a matrix equation, only B1 M1 A1 M1 A1.} \\
& {\(\begin{pmatrix}x\\y\\z\end{pmatrix} = \frac1{58}\begin{pmatrix}-72&-6&14\\150&-2&-5\\-20&8&-9\end{pmatrix}\begin{pmatrix}900\,000\\21\,000\,000\\15\,120\,000\end{pmatrix}\)} & dM1 & {Multiplies their constants on the left by the inverse of their coefficient matrix, by calculator or otherwise, to obtain values. Not for elimination alone.} \\
& {\(A\): \(3.6 \times 10^5\), \(B\): \(3.0 \times 10^5\), \(C\): \(2.4 \times 10^5\)} & A1 & {All three, each identified with its pass type; exact values accepted. Bare correct values after a correct matrix equation imply the dM1; values found by elimination alone, with no inverse used, score dM0 A0.} \\
\nopagebreak \SetCell[c=2]{r} Total for Question 6 & & 8 & \\
\end{msscheme}

\msQuestion{7}
\begin{msscheme}{
  hline{3-4,6-8} = {2-Z}{0.4pt, msGrid},
  hline{5,9} = {0.4pt, msGrid},
  row{9} = {font=\bfseries},
}
Question & Scheme & Marks & Notes \\
7(a) & {\(\det\mathbf{A} = a(3a-1) - 1(2a) = 3a^2 - 3a\)\\[4pt] \(3a(a-1) = 0 \Rightarrow a = \ldots\)} & M1 & {Expands \(\det\mathbf{A}\) by any row or column, or by Sarrus' rule, sets it equal to 0 and solves. Condone sign slips.} \\
& {\(a = 0\) or \(a = 1\)} & A1 & {Both values and no others.} \\
\nopagebreak & & [2] & \\
7(b) & {Minors: \(\begin{pmatrix} 3a-1 & 2a & 2 \\ a & a^2 & a \\ 1 & a & 3a-2 \end{pmatrix}\)} & M1 & {Forms a matrix of 2-by-2 determinants, each from deleting its entry's row and column. Condone slips in some entries.} \\
& {\(\operatorname{adj}\mathbf{A} = \begin{pmatrix} 3a-1 & -a & 1 \\ -2a & a^2 & -a \\ 2 & -a & 3a-2 \end{pmatrix}\)} & M1 & {Applies the alternating sign pattern and transposes their matrix of minors, in either order.} \\
& {\(\mathbf{A}^{-1} = \frac{1}{3a(a-1)}\begin{pmatrix} 3a-1 & -a & 1 \\ -2a & a^2 & -a \\ 2 & -a & 3a-2 \end{pmatrix}\)} & M1 A1 & {\msNote{M1}{Multiplies their adjugate, a 3-by-3 matrix, by the reciprocal of their quadratic determinant from (a).}\\ \msNote{A1}{Correct inverse, oe such as the factor taken inside the matrix or \(3a^2-3a\) unfactorised.}} \\
\nopagebreak & & [4] & \\
\nopagebreak \SetCell[c=2]{r} Total for Question 7 & & 6 & \\
\end{msscheme}

\msQuestion{8}
\begin{msscheme}{
  hline{3-4,6-7,9,11-12} = {2-Z}{0.4pt, msGrid},
  hline{5,8,10,13} = {0.4pt, msGrid},
  row{13} = {font=\bfseries},
}
Question & Scheme & Marks & Notes \\
8(a) & {\(\mathbf{AB}\): leading diagonal \(3k + 2 - 10 = 3k - 8\)\\[4pt] other entries \(0\), e.g. \(-4k + 2(k-4) + (2k+8) = 0\)\\[4pt] \(\Rightarrow \mathbf{AB} = (3k-8)\mathbf{I}\)} & M1 & {Forms at least one diagonal entry of \(\mathbf{AB}\) as rows of \(\mathbf{A}\) times columns of \(\mathbf{B}\). Condone slips; off-diagonal entries need not be shown.} \\
& {\(c = 8\)} & A1 & \\
\nopagebreak & & [2] & \\
8(b) & {\(\mathbf{A}^{-1} = \frac{1}{3k-8}\mathbf{B}\) unless\\[4pt] \(3k - \their{8} = 0 \Rightarrow k = \ldots\)} & M1 & {Sets their \(3k - c\) equal to zero, using part (a). Not for \(\det\mathbf{A} = 0\) alone.} \\
& {\(k = \frac{8}{3}\)} & A1 & \\
\nopagebreak & & [2] & \\
8(c) & {\(\mathbf{A}^{-1} = \frac{1}{3k-8}\mathbf{B}\)} & B1ft & {ft their \(c\): \(\frac{1}{3k-c}\mathbf{B}\), with \(\mathbf{B}\) named or written out.} \\
\nopagebreak & & [1] & \\
8(d) & {\(x = \frac{3 - 8 - 9}{3k-8}\)\\[4pt] \(y = \frac{1 + 2(k-4) - 3}{3k-8}\)\\[4pt] \(z = \frac{-10 + 2(2k+8) + 3(3k+2)}{3k-8}\)} & M1 & {Multiplies their \(\mathbf{A}^{-1}\) from (c), of the form \(\lambda\mathbf{B}\), by the column of constants \(1, 2, 3\), for at least one component. Condone slips.} \\
& {\(x = -\frac{14}{3k-8}\), \(y = \frac{2k-10}{3k-8}\)\\[4pt] \(z = \frac{13k+12}{3k-8}\)} & A1 A1 & {\msNote{A1}{Any two correct, in any form.}\\ \msNote{A1}{All three, each a single fraction with terms collected; equivalents such as \(\frac{14}{8-3k}\) or \(\frac{2(k-5)}{3k-8}\) accepted.}} \\
\nopagebreak & & [3] & \\
\nopagebreak \SetCell[c=2]{r} Total for Question 8 & & 8 & \\
\end{msscheme}

\msQuestion{9}
\begin{msscheme}{
  hline{3-4,6-10,12-14,16,18-19} = {2-Z}{0.4pt, msGrid},
  hline{5,11,15,17,20} = {0.4pt, msGrid},
  row{20} = {font=\bfseries},
}
Question & Scheme & Marks & Notes \\
9(a) & {\(\det(\mathbf{A}) = 1(1) + 2(4+a) + 3(0-a)\)} & M1 & {Expands along any row or column with three \(2 \times 2\) determinants, or uses the rule of Sarrus. Condone sign slips.} \\
& {\(\phantom{\det(\mathbf{A})} = 9 - a\)} & A1 & \\
\nopagebreak & & [2] & \\
9(b)(i) & {Minors: \(\begin{pmatrix}1 & a+4 & -a\\ -2 & 1-3a & 2a\\ -1 & -13 & 9\end{pmatrix}\)} & M1 & {Attempts the matrix of minors, each entry a \(2 \times 2\) determinant from the correct rows and columns. Condone arithmetic slips.} \\
& {\(\operatorname{adj}(\mathbf{A}) = \begin{pmatrix}1 & 2 & -1\\ -a-4 & 1-3a & 13\\ -a & -2a & 9\end{pmatrix}\)} & M1 & {Applies the alternating sign pattern to their minors and transposes, in either order.} \\
& {\(\mathbf{A}^{-1} = \frac{1}{9-a}\begin{pmatrix}1 & 2 & -1\\ -a-4 & 1-3a & 13\\ -a & -2a & 9\end{pmatrix}\)} & M1 A1 & {\msNote{M1}{Divides their adjugate by their determinant from (a).}\\ \msNote{A1}{Correct inverse oe, with \(\frac{1}{9-a}\) inside or outside the matrix.}} \\
& {\(\begin{pmatrix}x\\ y\\ z\end{pmatrix} = \mathbf{A}^{-1}\begin{pmatrix}-6\\ 7\\ a-1\end{pmatrix} = \frac{1}{9-a}\begin{pmatrix}9-a\\ 18-2a\\ a-9\end{pmatrix}\)} & dM1 & {Multiplies their \(\mathbf{A}^{-1}\) by the column of constants, inverse on the left, attempting all three entries.} \\
& {\(x = 1,\ y = 2,\ z = -1\)} & A1 & {cso, from the product with \(9-a\) cancelled.} \\
\nopagebreak & & [6] & \\
9(b)(ii) & {\(a = 9\): \((9, 0, 1) = \lambda(1, -2, 3) + \mu(4, 1, -1)\)\\[4pt] \(\phantom{a = 9:} \Rightarrow \lambda = \ldots,\ \mu = \ldots\)} & M1 & {Substitutes \(a = 9\) and seeks a combination of the first two equations giving the third, or eliminates a variable between pairs.} \\
& {Third \(=\) first \(+\ 2 \times\) second, with \(8 = -6 + 2(7)\)} & A1 & {Correct combination, checked for the constants as well as the coefficients, oe by elimination giving consistent equations.} \\
& {Sheaf: three distinct planes meeting in a common line} & E1 & {States a sheaf, or three planes meeting in one common line, after the A1. Infinitely many solutions alone is not enough.} \\
\nopagebreak & & [3] & \\
9(c)(i) & {\(9 - a < 0 \Rightarrow a > 9\)} & B1 & \\
\nopagebreak & & [1] & \\
9(c)(ii) & {\(12|9-a| < 12 \Rightarrow |9-a| < 1\)} & M1 & {Forms \(|\text{their } \det(\mathbf{A})| < 1\), or the equivalent for volumes, or considers both signs of the determinant.} \\
& {\(8 < a < 10\)} & A1 & {oe, such as \((8, 10)\). Condone \(a = 9\) excluded.} \\
\nopagebreak & & [2] & \\
\nopagebreak \SetCell[c=2]{r} Total for Question 9 & & 14 & \\
\end{msscheme}

\msQuestion{10}
\begin{msscheme}{
  hline{3-6,8-11} = {2-Z}{0.4pt, msGrid},
  hline{7,12} = {0.4pt, msGrid},
  row{12} = {font=\bfseries},
}
Question & Scheme & Marks & Notes \\
10(a) & {Unique point of intersection\\[4pt] \(\iff \det\begin{pmatrix}1&2&-1\\3&k&2\\0&k&4\end{pmatrix} \neq 0\)} & B1 & \\
& {\(\det = 1(4k-2k) - 2(12) - 1(3k)\)} & M1 & {Expands the determinant of the coefficient matrix by any valid method to an expression in \(k\). Condone sign slips.} \\
& {\(\phantom{\det} = -(k+24)\)} & A1 & {oe, such as \(-k-24\).} \\
& {Unique point for all \(k\) except \(k = -24\)} & A1 & {cso. Needs the conclusion that \(k = -24\) is the only exception.} \\
\nopagebreak & & [4] & \\
10(b) & {Second equation \(-\) third: \(3x - 2z = 7\)\\[4pt] \(3\times\) first \(-\) \((3x - 2z = 7)\): \(6y - z = 5\)} & M1 A1 & {\msNote{M1}{Combines the equations to eliminate \(k\), reaching two equations free of \(k\). Condone arithmetic slips.}\\ \msNote{A1}{Both correct, oe such as \(x = 4y - 1\) and \(z = 6y - 5\).}} \\
& {\(ky + 4(6y - 5) = k + 4 \Rightarrow (k+24)y = k+24\)\\[4pt] \(k \neq -24 \Rightarrow y = \ldots\)} & dM1 & {Substitutes their equations into one containing \(k\) to reach one equation in one unknown, and divides by its coefficient, using \(k \neq -24\).} \\
& {\(y = 1\)} & A1 & \\
& {\(x = 3,\ z = 1\): the point \((3, 1, 1)\), independent of \(k\)} & A1 & {cso. Both values, with the point stated to be independent of \(k\); A0 if the conclusion is missing.} \\
\nopagebreak & & [5] & \\
\nopagebreak \SetCell[c=2]{r} Total for Question 10 & & 9 & \\
\end{msscheme}

\end{document}
